\documentclass[DIV=calc]{scrartcl}
\usepackage[utf8]{inputenc}
\usepackage[T1]{fontenc}
\usepackage[ngerman]{babel}
\usepackage{graphicx}

\usepackage{fixltx2e}
\usepackage{ellipsis}
\usepackage[tracking=true]{microtype}

\usepackage{lmodern}                        % Ersatz fuer Computer Modern-Schriften
\usepackage{hfoldsty}

\usepackage{fourier}             % Schriftart
\usepackage[scaled=0.81]{helvet}     % Schriftart

\usepackage{url}
\usepackage{tocloft}             % Paket für Table of Contents

\usepackage{xcolor}
\definecolor{urlred}{HTML}{660000}

\usepackage{hyperref}
\hypersetup{
    colorlinks=true,    
    linkcolor=black,    % Farbe der internen Links (u.a. Table of Contents)
    urlcolor=black,    % Farbe der url-links
    citecolor=black} % Farbe der Literaturverzeichnis-Links

\usepackage{mdwlist}     % Änderung der Zeilenabstände bei itemize und enumerate

\parindent 0pt                 % Absatzeinrücken verhindern
\parskip 12pt                 % Absätze durch Lücke trennen

%\usepackage{titlesec}    % Abstand nach Überschriften neu definieren
%\titlespacing{\subsection}{0ex}{3ex}{-1ex}
%\titlespacing{\subsubsection}{0ex}{3ex}{-1ex}        

% \pagestyle{empty}
\setlength{\textheight}{23cm}
\usepackage{fancyhdr}
\pagestyle{fancy}
\cfoot{}
\lfoot{Zusammenkunft aller Physik-Fachschaften}
\rfoot{www.zapfev.de\\stapf@zapf.in}
\renewcommand{\headrulewidth}{0pt}
\renewcommand{\footrulewidth}{0.1pt}
\newcommand{\gen}{*innen}

\begin{document}
    \hspace{0.87\textwidth}
    \begin{minipage}{120pt}
        \vspace{-1.8cm}
        \includegraphics[width=80pt]{logo.pdf}
        \centering
        \small Zusammenkunft aller Physik-Fachschaften
    \end{minipage}
    \begin{center}
        \huge{Resolution zur Schaffung permanenter Stellen im wissenschaftlichen Mittelbau} \\
        \normalsize
    \end{center}
    
    \vspace{1cm}    
     Im wissenschaftlichen Mittelbau fehlt es an unbefristeten Stellen. Durch die drastische Reduktion solcher Stellen hat sich dieses Problem in den vergangenen Jahren noch erheblich verschärft\footnote{Der Anteil der unbefristeten Stellen am wissenschaftlichen und künstlerischen Personal sank zwischen 2004 und 2014 von 25,2\% auf 16,6\% (Statistisches Bundesamt, 2016: Hochschulen auf einen Blick)}.

Die Aussichten auf eine unbefristete Stelle in der Wissenschaft abseits der Professur sind derzeit verschwindend gering.
Durch die damit verbundene Unsicherheit sind Karriere und Familienplanung in der Wissenschaft in den seltensten Fällen vereinbar, was diesen Karriereweg für viele Studierende sehr unattraktiv macht.

Im Ergebnis wirkt sich dieses Umfeld prekärer Beschäftigungsverhältnisse zwangsläufig negativ auf die Kontinuität von Lehre und langfristigen Forschungsvorhaben aus.

Eine Beschleunigung des wissenschaftlichen Karrierewegs zur Professur ist der falsche Ansatz. Insbesondere punktuelle Initiativen, wie das aktuelle Vorhaben zur Förderung von Tenure-Track-Stellen, konzentrieren sich zu sehr auf die Professur als Karriereziel und verlieren die Schaffung einer ausgeglichenen akademischen Landschaft aus den Augen.

Nur durch eine deutliche Erhöhung der Anzahl unbefristeter Stellen im wissenschaftlichen Mittelbau kann es zu einer nachhaltigen Qualitätssicherung in der Forschung und Lehre, effizientem Wissenstransfer und einer Steigerung der Attraktivität der Karriere in der Wissenschaft kommen.

Daher fordert die ZaPF, auch im Interesse der Wissenschaft im deutschsprachigen Raum, die Schaffung unbefristeter Stellen im wissenschaftlichen Mittelbau.
    
    \vfill
    \begin{flushright}
        Verabschiedet am 27.05.2017 in Berlin
    \end{flushright}

    
\end{document}

%%% Local Variables:
%%% mode: latex
%%% TeX-master: t
%%% End:
